\documentclass{article}

% packages
\usepackage{amsmath, amsthm, thmtools, amsfonts, amssymb, luacode, catchfile, tikzducks, hyperref, ifthen}
\ifcsname c@kobocompile\endcsname
	\usepackage[a5paper, total={1072pt, 1448pt}, margin=10pt, includeheadfoot]{geometry} % set page margins
\else
	\usepackage[a4paper, margin=50pt, includeheadfoot]{geometry}
\fi
\usepackage[shortlabels]{enumitem}
\usepackage[skip=3pt, indent=0pt]{parskip}

% language
\usepackage[bidi=basic, layout=tabular, provide=*]{babel}
\ifcsname c@english\endcsname
	\babelprovide[main, import]{english}
\else
	\babelprovide[main, import]{hebrew}
	\babelprovide{rl}
\fi
%\babelfont{rm}{Libertinus Serif}
\babelfont{rm}[Renderer=Harfbuzz]{Libertinus Serif}
\babelfont{sf}{Libertinus Sans}
\babelfont{tt}{Libertinus Mono}

% style
\AddToHook{cmd/section/before}{\clearpage}	% Add line break before section
\linespread{1.3}
\setcounter{secnumdepth}{0}		% Remove default number tags from sections, this won't do well with theorems
\AtBeginDocument{\setlength{\belowdisplayskip}{3pt}}
\AtBeginDocument{\setlength{\abovedisplayskip}{3pt}}
\graphicspath{ {../images/} }

% operators
\DeclareMathOperator\cis{cis}
\DeclareMathOperator\Sp{Sp}
\DeclareMathOperator\tr{tr}
\DeclareMathOperator\im{Im}
\DeclareMathOperator\re{Re}
\DeclareMathOperator\diag{diag}
\DeclareMathOperator*\lowlim{\underline{lim}}
\DeclareMathOperator*\uplim{\overline{lim}}
\DeclareMathOperator\rng{rng}
\DeclareMathOperator\Sym{Sym}
\DeclareMathOperator\Arg{Arg}
\DeclareMathOperator\Log{Log}
\DeclareMathOperator\dom{dom}
\DeclareMathOperator\supp{Supp}
\DeclareMathOperator\var{Var}
\DeclareMathOperator\cov{Cov}

% commands
%\renewcommand\qedsymbol{\textbf{מש''ל}}
%\renewcommand\qedsymbol{\fbox{\emoji{lizard}}}
\newcommand{\Aa}[0]{\mathcal{A}}
\newcommand{\Bb}[0]{\mathcal{B}}
\newcommand{\CC}[0]{\mathbb{C}}
\newcommand{\Cc}[0]{\mathcal{C}}
\newcommand{\EE}[0]{\mathbb{E}}
\newcommand{\FF}[0]{\mathbb{F}}
\newcommand{\Ff}[0]{\mathcal{F}}
\newcommand{\Ii}[0]{\mathcal{I}}
\newcommand{\Gg}[0]{\mathcal{G}}
\newcommand{\Ll}[0]{\mathcal{L}}
\newcommand{\Mm}[0]{\mathcal{M}}
\newcommand{\NN}[0]{\mathbb{N}}
\newcommand{\Nn}[0]{\mathcal{N}}
\newcommand{\PP}[0]{\mathbb{P}}
\newcommand{\Pp}[0]{\mathcal{P}}
\newcommand{\QQ}[0]{\mathbb{Q}}
\newcommand{\RR}[0]{\mathbb{R}}
\newcommand{\Rr}[0]{\mathcal{R}}
\newcommand{\Ss}[0]{\mathcal{S}}
\newcommand{\TT}[0]{\mathbb{T}}
\newcommand{\Uu}[0]{\mathcal{U}}
\newcommand{\Vv}[0]{\mathcal{V}}
\newcommand{\Ww}[0]{\mathcal{W}}
\newcommand{\ZZ}[0]{\mathbb{Z}}
\newcommand{\acts}[0]{\circlearrowright}
\newcommand{\explain}[2] {
	\begin{flalign*}
		 && \text{#2} && \text{#1}
	\end{flalign*}
}
\newcommand{\maketitleprint}[0]{ \begin{center}
	%\begin{tikzpicture}[scale=3]
	%	\duck[graduate=gray!20!black, tassel=red!70!black]
	%\end{tikzpicture}	
	\includegraphics[width=6cm]{cover}
\end{center}
}

% theorem commands
\newtheoremstyle{c_remark}
	{}	% Space above
	{}	% Space below
	{}% Body font
	{}	% Indent amount
	{\bfseries}	% Theorem head font
	{}	% Punctuation after theorem head
	{.5em}	% Space after theorem head
	{\thmname{#1}\thmnumber{ #2}\thmnote{ \normalfont{\text{(#3)}}}}	% head content
\newtheoremstyle{c_definition}
	{3pt}	% Space above
	{3pt}	% Space below
	{}% Body font
	{}	% Indent amount
	{\bfseries}	% Theorem head font
	{}	% Punctuation after theorem head
	{.5em}	% Space after theorem head
	{\thmname{#1}\thmnumber{ #2}\thmnote{ \normalfont{\text{(#3)}}}}	% head content
\newtheoremstyle{c_plain}
	{3pt}	% Space above
	{3pt}	% Space below
	{\itshape}% Body font
	{}	% Indent amount
	{\bfseries}	% Theorem head font
	{}	% Punctuation after theorem head
	{.5em}	% Space after theorem head
	{\thmname{#1}\thmnumber{ #2}\thmnote{ \text{(#3)}}}	% head content

\ifcsname c@english\endcsname
	\theoremstyle{plain}
	\newtheorem{theorem}{Theorem}[section]
	\newtheorem{lemma}[theorem]{Lemma}
	\newtheorem{proposition}[theorem]{Proposition}
	\newtheorem*{proposition*}{Proposition}
	%\newtheorem{corollary}[theorem]{אין חלופה עברית}

	\theoremstyle{definition}
	\newtheorem{definition}[theorem]{Definition}
	\newtheorem*{definition*}{Definition}
	\newtheorem{example}{Example}[section]
	\newtheorem{exercise}{Exercise}[section]

	\theoremstyle{remark}
	\newtheorem*{remark}{Remark}
	\newtheorem*{solution}{Solution}
	\newtheorem{conclusion}[theorem]{Conclusion}
	\newtheorem{notation}[theorem]{Notation}
\else
	\theoremstyle{c_plain}
	\newtheorem{theorem}{משפט}[section]
	\newtheorem{lemma}[theorem]{למה}
	\newtheorem{proposition}[theorem]{טענה}
	\newtheorem*{proposition*}{טענה}
	%\newtheorem{corollary}[theorem]{אין חלופה עברית}

	\theoremstyle{c_definition}
	\newtheorem{definition}[theorem]{הגדרה}
	\newtheorem*{definition*}{הגדרה}
	\newtheorem{example}{דוגמה}[section]
	\newtheorem{exercise}{תרגיל}[section]

	\theoremstyle{c_remark}
	\newtheorem*{remark}{הערה}
	\newtheorem*{solution}{פתרון}
	\newtheorem{conclusion}[theorem]{מסקנה}
	\newtheorem{notation}[theorem]{סימון}
\fi

% Questions related commands
\newcounter{question}
\setcounter{question}{1}
\newcounter{sub_question}
\setcounter{sub_question}{1}

\ifcsname c@english\endcsname
	\newcommand{\question}[1][0]{
		\ifthenelse{#1 = 0}{}{\setcounter{question}{#1}}
		\section{Question \arabic{question}}
		\addtocounter{question}{1}
		\setcounter{sub_question}{1}
	}

	\newcommand{\subquestion}[1][0]{
		\ifthenelse{#1 = 0}{}{\setcounter{sub_question}{#1}}
		\subsection{Part \alph{sub_question}}
		\addtocounter{sub_question}{1}
	}
\else
	\newcommand{\question}[1][0]{
		\ifthenelse{#1 = 0}{}{\setcounter{question}{#1}}
		\section{שאלה \arabic{question}}
		\addtocounter{question}{1}
		\setcounter{sub_question}{1}
	}

	\newcommand{\subquestion}[1][0]{
		\ifthenelse{#1 = 0}{}{\setcounter{sub_question}{#1}}
		\subsection{סעיף \localecounter{letters.gershayim}{sub_question}}
		\addtocounter{sub_question}{1}
	}
\fi

% import lua and start of document
\directlua{common = require ('../common')}

\GetEnv{AUTHOR}

% headers
\author{\AUTHOR}
\date\today

\title{פתרון מטלה 06 --- אנליזה על יריעות, 80426}

\DeclareMathOperator{\vol}{vol}

\begin{document}
\maketitle
\maketitleprint[blue]

\question{}
תהי $U \subseteq \RR^k$ פתוחה ותהי $\varphi : U \to \RR^n$ כך ש־$(\varphi(U), \varphi)$ היא יריעה פרמטרית רגולרית $k$־מימדית.
נראה שלכל $x_0 \in U$ קיימת סביבה $x_0 \in U' \subseteq U$ כך ש־$\varphi(U')$ היא יריעה.
\begin{proof}
	יהי $x_0 \in U$ ונבחין כי $\deg(D \varphi \mid_{x_0}) = k$, לכן תנאי משפט הפונקציה ההפוכה חלים,
	זאת שכן נוכל לבחון את ההעתקה לתוך $\varphi(U)$ תוך שימוש בהעתקה אורתוגונלית מתאימה,
	וקיימת סביבה $x_0 \in U' \subseteq U$ בה הפונקציה $\varphi$ היא הפיכה.
	נגדיר $\psi = \varphi \restriction U'$, אז $\psi$ הפיכה, וכן מהרחבה למשפט הפונקציה ההפוכה היא גם גזירה, כלומר $\varphi$ היא דיפאומורפיזם, בפרט $\psi$ היא הומיאומורפיזם,
	כלומר היא הפיכה, גזירה, ופתוחה, וכמובן רגולרית כצמצום של העתקה רגולרית.

	נגדיר את $M = \psi(U')$, ונטען כי זוהי יריעה $k$־מימדית.
	לכל נקודה $x_1 \in M$, נבחין כי נוכל לבחור את $\psi$ ולקבל את התנאים הדרושים להגדרת יריעה, ולכן $M$ אכן יריעה.
\end{proof}

\question{}
תהי $M \subseteq \RR^n$ יריעה $k$־מימדית,
ותהי $\alpha : U \to W$ פרמטריזציה מקומית סביב $p \in M$. \\
נראה שקיימת קבוצה פתוחה $\tilde{W} \subseteq \RR^n$ ו־$\psi : \tilde{W} \to U$ חלקה,
כך ש־$p \in \tilde{W}$ וגם,
\[
	\forall q \in \alpha^{-1}(\tilde{W}),\ 
	\psi(\alpha(q)) = q
\]
\begin{proof}
	נניח ש־$\alpha(x_0) = p$ עבור $x_0 \in U$.
	נניח ללא הגבלת הכלליות ש־$k$ השורות הראשונות של $D \varphi \mid_{x_0}$ בלתי־תלויות לינארית,
	מותר לנו להניח כן שכן אחרת נוכל לכפול בהעתקה אורתוגונלית מתאימה, וכי נתון כי $M$ יריעה ולכן ההעתקה $\alpha$ היא רגולרית.
	נגדיר את הקבוצה $\tilde{U} = U \times \RR^{n - k}$ וכן נגדיר את ההעתקה $\tilde{\alpha} : \tilde{U} \to \RR^n$ על־ידי,
	\[
		\tilde{\alpha}(x_1, \ldots, x_n)
		= \alpha(x_1, \ldots, x_k) + \sum_{i = k + 1}^n x_i \cdot e_i
	\]
	עבור $e_i$ איבר הבסיס הסטנדרטי עבור $1 \le i \le n$.
	מבדיקה ישירה מתקיים,
	\[
		D \tilde{\varphi} \mid_{(x_0, 0)}
		= \begin{pmatrix}
			D \varphi \mid_{x_0} & 0 \\
			\vdots & \id
		\end{pmatrix}
	\]
	כלומר זוהי מטריצה כך שהיא מרחיבה את $D \varphi \mid_{x_0}$ על־ידי מטריצת היחידה.
	נבחין כי החישוב מתקבל ישירות מההגדרה בהתאם למימד.
	אנו יודעים כי דרגת המטריצה $k + (n - k) = n$, כלומר זוהי מטריצה רגולרית, ולכן הנגזרת היא בעלת דטרמיננטה לא אפס, ונוכל להסיק שתנאי משפט הפונקציה ההפוכה חלים.
	נסיק כי קיימת $U_0 \times V_0 \subseteq \tilde{U}$ כך ש־$x_0 \in U_0 \subseteq U$, וקיימת $p \in \tilde{W} \subseteq \RR^n$, כך ש־$\tilde{\alpha} : U_0 \times V_0 \to \tilde{W}$ דיפאומורפיזם.
	נגדיר את $\psi : \tilde{W} \to U$ על־ידי $\psi = \pi_k \circ \tilde{\alpha}^{-1}$ עבור $\pi_k : \RR^n \to \RR^k$ ההטלה של $k$ האיברים הראשונים.
	נבחין כי ממשפט ההעתקה ההפוכה $\psi$ היא העתקה גזירה, ומהפעלה חוזרת ונשנית נקבל שהיא גזירה מכל סדר, קרי היא פונקציה חלקה.
	לבסוף נובע שעבור $q \in U_0 = \alpha^{-1}(\tilde{W})$ מתקיים $\psi(\alpha(q)) = q$.
\end{proof}

\question{}
תהי $M$ יריעה $k$־מימדית ב־$\RR^n$, נגדיר את האגד המשיק,
\[
	TM
	= \{ (p, v) \in \RR^n \times \RR^n \mid p \in M, v \in T_p M \}
\]

\subquestion{}
נראה ש־$TM$ היא יריעה $2k$־מימדית.
\begin{proof}
	תהי נקודה $p \in M$ כלשהי, נבחין כי $M$ יריעה ולכן קיימת פרמטריזציה מקומית $\alpha : U \to M$ כך ש־$\alpha(0) = p$ (ללא הגבלת הכלליות).
	נגדיר את הפונקציה $\overline{\alpha} : U \times \RR^k \to \RR^{2n}$ על־ידי $\overline{\alpha}(q, v) = (\alpha(q), d \alpha \mid_p(v))$.
	לכל $u \in T_p M$ נבחין כי קיים $v \in \RR^k$ כך ש־$u = T_p(M) \cdot v$, ולכן $\overline{\alpha}(0, v) = (p, u)$, כלומר $\overline{\alpha}$ חשודה כפרמטריזציה של $(p, u)$.

	אנו יודעים כי $\alpha$ הפיכה, וכן $d \alpha \mid_p$ היא העתקה לינארית רגולרית ולכן הפיכה אף היא ב־$T_p M$, ולכן נסיק שגם $\overline{\alpha}$ היא הפיכה.
	אנו גם יודעים כי $\alpha$ חלקה ורגולרית, וכן כל העתקה לינארית היא חלקה, ומרגולריות $\alpha$ נסיק שנגזרתה רגולרית אף היא בתחום, ולכן נסיק שגם $\overline{\alpha}$ היא חלקה ורגולרית.
	עלינו להראות שההעתקה היא פתוחה, אך גם הפעם, $\alpha$ היא העתקה פתוחה, והעתקות לינאריות תמיד פתוחות, ולכן נסיק ש־$\overline{\alpha}$ היא אכן פרמטריזציה מקומית של $(p, u)$.

	נסיק אם כן ש־$TM$ היא יריעה $2k$־מימדית ב־$\RR^{2n}$.
\end{proof}

\subquestion{}
נראה שאם $M$ היא תת־קבוצה פתוחה של $\RR^n$, אז $TM = M \times \RR^n$.
\begin{proof}
	נבהיר תחילה שכמסקנה מהטענה מהתרגול $M$ היא יריעה $n$־מימדית, ולכן $TM$ יריעה $2n$־מימדית ב־$\RR^{2n}$.
	לכל $p \in M$ נניח ש־$\alpha : U \to M$ פרמטריזציה מקומית, אז $\alpha$ רגולרית, דהינו $\dim D \alpha = n$ ולכן היא הפיכה (ממשפט הפונקציה ההפוכה), ובפרט לכל $u \in \RR^n$ נסיק ש־$u \in T_p M$.
	נובע אם כך ש־$\{ p \} \times \RR^n \in TM$ לכל $p \in M$, ונסיק ש־$TM = M \times \RR^n$.
\end{proof}

\question{}
תהיינה $M \subseteq \RR^m$ ו־$N \subseteq \RR^n$ שתי יריעות.
תהי $f : M \to N$.
נניח ש־$\alpha : U \to W$ פרמטריזציה מקומית של $M$ סביב $p \in M$,
ונניח ש־$\beta : A \to B$ פרמטריזציה מקומית של $N$ סביב $f(p)$. \\
נראה שאם קיימת סביבה פתוחה $U' \subseteq U$ של $\alpha^{-1}(p)$ כך ש־$\beta^{-1} \circ f \circ \alpha : U' \to A$ היא חלקה, אז $f$ חלקה ב־$p$.
\begin{proof}
	נניח שקיימת סביבה פתוחה $U'$ כך שהתנאי מתקיים.
	נבחין תחילה כי $\beta^{-1}(f(\alpha(x_0))) = \beta^{-1}(f(p))$ היא נקודה פנימית בהרכבה.
	אנו גם יודעים שהרכבת פונקציות חלקות היא פונקציה חלקה, לכן גם $\beta \circ \beta^{-1} \circ f \circ \alpha = f \circ \alpha : U' \to N$ היא העתקה חלקה ב־$\alpha^{-1}(p)$.
	מאותה סיבה בדיוק גם $f \circ \alpha \circ \alpha^{-1} = f : M \to N$ היא חלקה ב־$\alpha^{-1}(\alpha(x_0)) = p$ ונקבל ש־$f$ חלקה ב־$p$, כפי שרצינו.
	נבהיר ונאמר ש־$\alpha, \beta, \alpha^{-1}, \beta^{-1}$ כולן דיפאומורפיזמים חלקים, ולכן הפיכות, חלקות, ומוגדרות בסביבות הנתונות.
\end{proof}

\end{document}
